\documentclass[%
 reprint,
 amsmath,amssymb,
 aps,
 prd,nofootinbib,nobibnotes,
 superscriptaddress,groupedaddress,
]{revtex4-2}

\bibliographystyle{apsrev4-1}

\usepackage{graphicx}
\usepackage{dcolumn}
\usepackage{bm}
\usepackage[dvipsnames, usenames]{xcolor}
\usepackage{acronym}
\usepackage[hidelinks]{hyperref}
\usepackage{physics}
\usepackage{orcidlink}

\graphicspath{{referee-plots}}

\newcommand{\red}[1]{{\color{Red} #1}}
\newcommand{\jh}[1]{{\color{orange}JH: #1}}
\newcommand{\mm}[1]{{\color{cyan}MM: #1}}
\newcommand{\sa}[1]{{\color{violet}SA: #1}}
\newcommand{\sv}[1]{{\color{BurntOrange}SV: #1}}
\newcommand{\comment}[1]{}

\newcommand{\ligo}{\affiliation{LIGO Laboratory, Massachusetts Institute of Technology, Cambridge, MA 02139, USA}}
\renewcommand{\mit}{\affiliation{Kavli Institute for Astrophysics and Space Research and Department of Physics, Massachusetts Institute of Technology, Cambridge, MA 02139, USA}}

\begin{document}

\begin{figure*}
\centering
\includegraphics[width=0.49\linewidth]{chiq_double_correlated.pdf}
\includegraphics[width=0.49\linewidth]{chiq_double_correlated_z0is1.0.pdf}
\caption{\textit{Left.} Inferred comoving merger rate density $\mathcal{R}(\chi_\mathrm{eff}, q;z=0.2)$ at a fixed redshift $z=0.2$ for a simulated population in which there is a three-dimensional correlation between mass ratio, redshift and $\chi_\mathrm{eff}$. \textsc{PixelPop} is used to measure the merger rate jointly over mass ratio $q$ and effective spin $\chi_\mathrm{eff}$, but neglects the correlation with redshift. The central panel displays the two-dimensional posterior median, with lower to higher values shaded darker to brighter. The underlying true population is also sliced at $z=0.2$, and visualized with dashed black lines enclosing 50\%, 90\%, and 99\% of the distribution. The upper panel shows slices $\mathcal{R}(\chi_\mathrm{eff}=0,q)$ (blue) and $\mathcal{R}(\chi_\mathrm{eff}=0.25,q)$ (orange). The solid lines gives the posterior median and the shaded regions encloses the 90\% posterior credible regions. The dashed lines give the marginal distributions of the true population. Similarly, the right-hand panel shows slices $\mathcal{R}(\chi_\mathrm{eff},q=1)$ (blue) and $\mathcal{R}(\chi_\mathrm{eff},q=0.6)$ (orange).\textit{Right.} The same as on the left hand size, with the redshift sliced at $z=1.0$.}
\label{fig:qchi_doublecor}
\end{figure*}


\begin{figure*}
\centering
\includegraphics[width=0.49\linewidth]{spearman_correlations.pdf}
\includegraphics[width=0.49\linewidth]{spearman_correlations_z0is1.0.pdf}
\caption{Posterior (solid lines) and prior (dashed lines) distributions for the spearman correlation coefficient for the synthetic population with a three dimensional correlation, and when PixelPop only models the correlation between mass ratio and effective spin. On the left hand panel is when the redshift is sliced at $z=0.2$, and on the right is when the redshift is sliced at $z=1.0$. The true value at these redshift slices is marked with a vertical black line. The histograms look the same because the slicing at different redshifts represents a constant rescaling for each hyperposterior sample--this is baked in the assumption that redshift is uncorrelated with mass ratio and effective spins. The Spearman correlation is invariant to a constant rescaling.
}
\label{fig:Spearman_zchi_uncor}
\end{figure*}

\begin{figure*}
\centering
\includegraphics[width=0.49\linewidth]{chiz_double_correlated.pdf}
\includegraphics[width=0.49\linewidth]{chiz_double_correlated_q0is0.6.pdf}
\caption{\textit{Left.} Inferred comoving merger rate density $\mathcal{R}(\chi_\mathrm{eff},q=1;z)$ for a simulated population in which there is a three-dimensional correlation between mass ratio, redshift and $\chi_\mathrm{eff}$. \textsc{PixelPop} is used to measure the merger rate jointly over redshift $z$ and effective spin $\chi_\mathrm{eff}$, but neglects the correlation with mass ratio. The central panel displays the two-dimensional posterior median, with lower to higher values shaded darker to brighter. The underlying true population is also sliced at $q=1$ equal mass ratios, and is visualized with dashed black lines enclosing 50\%, 90\%, and 99\% of the distribution. The upper panel shows slices $\mathcal{R}(\chi_\mathrm{eff}=0;z)$ (blue) and $\mathcal{R}(\chi_\mathrm{eff}=0.25;z)$ (orange). The solid lines gives the posterior median and the shaded regions encloses the 90\% posterior credible regions. The dashed lines give the marginal distributions of the true population. Similarly, the right-hand panel shows slices $\mathcal{R}(\chi_\mathrm{eff};z=0.1)$ (blue) and $\mathcal{R}(\chi_\mathrm{eff};z=0.75)$ (orange). \textit{Right.} The same as on the left hand side, where the mass ratio is fixed to the $q=0.6$ slice.}
\label{fig:zchi_doublecor}
\end{figure*}


\begin{figure*}
\centering
\includegraphics[width=0.49\linewidth]{spearman_broadening.pdf}
\includegraphics[width=0.49\linewidth]{spearman_broadening_q0is0.6.pdf}
\caption{Posterior (solid lines) and prior (dashed lines) distributions for the Spearman broadening coefficient for the synthetic population with a three dimensional correlation, and when PixelPop only models the correlation between redshift and effective spin. On the left is the Spearman broadening for the slice $q=1$, while on the right is the Spearman broadening coefficient for the slice $q=0.6$. The true value for each mass ratio slice is marked with a vertical black line. The histograms look the same because the slicing at different mass ratios represents a constant rescaling for each hyperposterior sample--this is baked in the assumption that mass ratio is uncorrelated with redshift and effective spins. The Spearman correlation is invariant to a constant rescaling.}
\label{fig:Spearman_zchi_doublecor}
\end{figure*}


\end{document}